\documentclass[a4paper,10pt]{article}
\usepackage[utf8]{inputenc}
\usepackage[T1]{fontenc}
\usepackage[spanish]{babel}
\usepackage{geometry}
\usepackage{array}
\usepackage{colortbl}
\usepackage{xcolor}
\usepackage{booktabs}
\usepackage{enumitem}
\usepackage{fancyhdr}
\usepackage{tikz}
\usepackage{tcolorbox}
\usepackage{longtable}
\usepackage{multicol}
\usepackage{amssymb}

\tcbuselibrary{skins,breakable}

\geometry{left=1.8cm, right=1.8cm, top=2.5cm, bottom=2cm}

% Colores principales
\definecolor{naranja}{HTML}{E67E22}
\definecolor{gris}{HTML}{F2F2F2}
\definecolor{grismedio}{HTML}{CCCCCC}
% Colores para interrogativos (gramática visual)
\definecolor{verde}{HTML}{27AE60}
\definecolor{colornaranja}{HTML}{E67E22}
\definecolor{rojo}{HTML}{E74C3C}
\definecolor{morado}{HTML}{8E44AD}
\definecolor{azul}{HTML}{2980B9}
\definecolor{rosa}{HTML}{E91E8B}
\definecolor{azulmasc}{HTML}{2980B9}
\definecolor{rosafem}{HTML}{E91E8B}
\definecolor{negro}{HTML}{000000}

% Header/footer
\pagestyle{fancy}
\fancyhf{}
\fancyhead[L]{\small\textbf{Nuevo Compañeros 1 — U03: La Familia}}
\fancyhead[R]{\small\textbf{Píldora formativa 3.3 — Interrogativos}}
\fancyfoot[C]{\thepage}

\setlength{\parindent}{0pt}
\renewcommand{\arraystretch}{1.3}

% Estilo de caja para secciones MARS
\newtcolorbox{marsbox}[1][]{
  colback=naranja!5,
  colframe=naranja,
  fonttitle=\bfseries\small,
  title=#1,
  boxrule=0.8pt,
  arc=2pt,
  left=4pt, right=4pt, top=2pt, bottom=2pt
}

% Estilo de caja para awareness / notas importantes
\newtcolorbox{notabox}[1][]{
  colback=gris,
  colframe=grismedio,
  fonttitle=\bfseries\small,
  title=#1,
  boxrule=0.6pt,
  arc=2pt,
  left=4pt, right=4pt, top=2pt, bottom=2pt
}

% Estilo de barra de diapositiva
\newcommand{\diapheader}[2]{%
  \noindent\begin{tcolorbox}[
    colback=naranja,
    colframe=naranja,
    coltext=white,
    fontupper=\bfseries\large,
    boxrule=0pt,
    arc=3pt,
    left=6pt, right=6pt, top=4pt, bottom=4pt
  ]
  DIAPOSITIVA #1 — #2
  \end{tcolorbox}
}

\begin{document}

%% ============================================================
%% PÁGINA 1: PORTADA + CONTENIDO PARA EL PROFESOR
%% ============================================================

\begin{center}
{\LARGE\textbf{\textcolor{naranja}{PÍLDORA FORMATIVA 3.3}}}\\[0.3cm]
{\Large\textbf{INTERROGATIVOS: SISTEMA DE 7 PARTÍCULAS}}\\[0.1cm]
{\Large\textbf{Y CONCORDANCIA DE CUÁNTO}}\\[0.3cm]
{\normalsize Sección Gramática (p.36--37) — Proyectable — 4 diapositivas}\\[0.1cm]
{\small Enfoque inductivo MARS EARS (Conti, 2023)}
\end{center}

\vspace{0.4cm}

%% --- Resumen de diapositivas ---
\begin{tcolorbox}[
  colback=naranja!8,
  colframe=naranja,
  title={\textbf{Resumen de diapositivas}},
  fonttitle=\bfseries,
  boxrule=0.8pt,
  arc=3pt
]
\renewcommand{\arraystretch}{1.4}
\begin{tabular}{c l l}
\textbf{Diap.} & \textbf{Título} & \textbf{Fase MARS} \\
\midrule
1 & Recordad estas preguntas & M (Modelling) \\
2 & ¿Cómo preguntamos por esta información? & A (Awareness Raising) \\
3 & ¿Cuántos o cuántas? ¿Quién o quiénes? & A (Awareness --- concordancia) \\
4 & ¡Gritad la pregunta! + ¿Con tilde o sin tilde? & S (Structured Production) \\
\end{tabular}
\end{tcolorbox}

\vspace{0.4cm}

%% --- Tabla de 7 interrogativos ---
\subsection*{\textcolor{naranja}{Sistema de 7 partículas interrogativas}}

Los interrogativos del español forman un sistema cerrado de 7 partículas, cada una asociada a un tipo de información:

\vspace{0.2cm}

\begin{center}
\renewcommand{\arraystretch}{1.4}
\begin{tabular}{|>{\bfseries}l|l|l|}
\hline
\rowcolor{naranja!15}
\textbf{Interrogativo} & \textbf{Pregunta por\ldots} & \textbf{Ejemplo} \\
\hline
Qué & cosa, actividad & ¿Qué estudias? \\
\hline
Quién / Quiénes & persona(s) & ¿Quién es esta chica? \\
\hline
Cuál & selección de un conjunto & ¿Cuál es tu asignatura favorita? \\
\hline
Dónde & lugar & ¿Dónde viven tus abuelos? \\
\hline
Cómo & manera, nombre & ¿Cómo se llama tu abuela? \\
\hline
Cuándo & tiempo & ¿Cuándo tenemos el examen? \\
\hline
Cuánto/-a/-os/-as & cantidad & ¿Cuántos hermanos tienes? \\
\hline
\end{tabular}
\end{center}

\vspace{0.3cm}

%% --- Dificultad principal ---
\begin{notabox}[Dificultad principal para A1.1: \textit{qué} vs. \textit{cuál}]
En español, \textbf{cuál} selecciona de un conjunto cerrado (\textit{¿Cuál es tu favorito?}) mientras que \textbf{qué} pregunta por una categoría abierta (\textit{¿Qué estudias?}). En inglés, \textit{what} cubre ambos usos, lo que genera interferencia.
\end{notabox}

\vspace{0.3cm}

%% --- Concordancia de cuánto ---
\subsection*{\textcolor{naranja}{Concordancia de \textit{cuánto}}}

\textbf{Cuánto} es el único interrogativo que varía morfológicamente --- tiene concordancia de género y número:

\vspace{0.2cm}

\begin{center}
\renewcommand{\arraystretch}{1.3}
\begin{tabular}{|l|l|l|}
\hline
\rowcolor{naranja!15}
\textbf{Forma} & \textbf{Tipo} & \textbf{Ejemplo} \\
\hline
cuánto & M sing. no contable & ¿Cuánto café bebes? \\
\hline
cuánta & F sing. no contable & ¿Cuánta agua bebes? \\
\hline
cuántos & M plural contable & ¿Cuántos hermanos tienes? \\
\hline
cuántas & F plural contable & ¿Cuántas hermanas tienes? \\
\hline
\end{tabular}
\end{center}

\vspace{0.2cm}

Todos los interrogativos llevan tilde. Este dato visual ayuda al estudiante a identificarlos en textos escritos.

\vspace{0.3cm}

%% --- Lo que el estudiante ya trae ---
\begin{notabox}[Lo que el estudiante ya trae]
Ha usado interrogativos informalmente en U01--U02 y en Vocabulario (\textit{¿Cómo se llama?, ¿Quién es?, ¿Dónde vive?}) y acaba de completar preguntas con verbos regulares (Bloque 1, Actividad 2). Los interrogativos aparecieron pero no se han formalizado como sistema.
\end{notabox}

\newpage

%% ============================================================
%% PÁGINA 2: TABLA MARS EARS
%% ============================================================

\begin{center}
{\Large\textbf{\textcolor{naranja}{MARS EARS — Correspondencia con Bloque 2}}}
\end{center}

\vspace{0.3cm}

{\small\textit{La M (Modelling) ocurrió parcialmente en Bloque 1 (Actividad 2: el estudiante completó preguntas con interrogativos sin formalizarlos) y en U01--U02. Ciclo completo inductivo: el patrón es deducible (cada interrogativo se asocia a un tipo de respuesta).}}

\vspace{0.4cm}

\renewcommand{\arraystretch}{1.5}
{\small
\begin{longtable}{|>{\bfseries\raggedright\arraybackslash}p{2.2cm}|p{4.8cm}|p{3.8cm}|p{3.2cm}|}
\hline
\rowcolor{naranja!15}
\textbf{MARS EARS} & \textbf{Qué hace el estudiante} & \textbf{Material / recurso} & \textbf{Fase} \\
\hline
\endfirsthead
\hline
\rowcolor{naranja!15}
\textbf{MARS EARS} & \textbf{Qué hace el estudiante} & \textbf{Material / recurso} & \textbf{Fase} \\
\hline
\endhead

M --- Modelling &
\textit{(parcialmente ocurrido)} En Bloque 1 Act.~2 completó preguntas con \textit{dónde, qué, cómo}. En U01--U02 usó \textit{¿Quién es?, ¿Cómo se llama?} &
Actividad 2 p.36 + U01--U02 &
Bloque 1 Fase 4 \\
\hline

A --- Awareness Raising &
Ve 4 preguntas-respuesta en la píldora. Nota que cada interrogativo se asocia a un tipo de respuesta (lugar$\to$dónde, persona$\to$quién). Categoriza los 7 interrogativos con tarjetas. Verifica con el cuadro del libro. &
Píldora proyectable (diap.~1--4) + tarjetas + cuadro p.36 &
Fase 5 \\
\hline

R --- Receptive Processing &
Selecciona el interrogativo correcto entre dos opciones (Actividad 3: 7 ítems de selección binaria). Reconoce sin producir libremente. &
Libro p.36 act.~3 &
Fase 6 \\
\hline

S --- Structured Production &
Completa huecos con \textit{cuánto/-a/-os/-as} (Actividad 4: concordancia). Weaning off: con cuadro $\to$ sin cuadro. &
Libro p.37 act.~4 &
Fase 7 \\
\hline

E --- Expansion &
Integra interrogativos + posesivos en Bloque 3 (Actividades 6--8). Repertorio cruzado. &
Libro p.37 act.~6--8 &
Bloque 3 \\
\hline

A --- Autonomous use &
Formula preguntas propias a compañeros sobre sus familias. &
Interacción oral &
Bloque 3 Fases 11--12 \\
\hline

R --- Routinisation &
Consolidación distribuida: tarea 24\,h, mención 1 semana, integrador 4 semanas. &
Cuaderno + clases futuras &
Cierre de sección \\
\hline

S --- Spontaneity &
Uso espontáneo en Comunicación y Destrezas. &
--- &
Secciones §3.3--§3.5 \\
\hline
\end{longtable}
}

\newpage

%% ============================================================
%% DIAPOSITIVA 1 — RECORDAD ESTAS PREGUNTAS
%% ============================================================

\diapheader{1}{RECORDAD ESTAS PREGUNTAS}

\vspace{0.2cm}

\begin{marsbox}[Fase MARS: M (Modelling)]
\textit{Técnica: Input flooding contextualizado con gramática visual (Conti, 2023; Romo Simón, 2014). El estudiante recibe 7 exposiciones a interrogativos dentro de un contexto que ya conoce (textos de Javier y Lucía, p.35). Los textos aparecen en pantalla como contexto visual. Debajo de cada texto, las preguntas-respuesta con el interrogativo y la parte de la respuesta que le corresponde marcados en el mismo color. Un icono al lado de cada pregunta indica el tipo de información. El estudiante percibe el vínculo interrogativo $\leftrightarrow$ respuesta de forma visual sin explicación.}
\end{marsbox}

\vspace{0.2cm}

%% --- Leyenda de iconos ---
\begin{center}
{\small
\fbox{\parbox{14cm}{\centering
\textcolor{verde}{\textbf{lugar}} \hspace{0.8cm}
\textcolor{colornaranja}{\textbf{nombre}} \hspace{0.8cm}
\textcolor{rojo}{\textbf{persona}} \hspace{0.8cm}
\textcolor{morado}{\textbf{cantidad}} \hspace{0.8cm}
\textcolor{azul}{\textbf{cosa / actividad}}
}}}
\end{center}

\vspace{0.3cm}

%% --- TEXTO 1: Javier ---
\subsection*{\textcolor{naranja}{Texto 1 --- Javier} {\small (visible en pantalla, el estudiante reconoce el texto ya trabajado)}}

{\small
``Javier Rodríguez vive \textcolor{verde}{\textbf{en Getafe}}, una ciudad muy cerca de Madrid. Javier tiene una hermana pequeña, se llama \textcolor{colornaranja}{\textbf{Alejandra}} y tiene 9 años, y en la misma casa vive también \textcolor{rojo}{\textbf{Manolo, su abuelo}}. El padre de Javier es camionero y de lunes a viernes trabaja fuera de Madrid. La madre, Catalina, es enfermera. Javier estudia el segundo curso de Educación Secundaria Obligatoria en el instituto de su barrio y juega al fútbol tres días a la semana.''}

\vspace{0.2cm}

\textit{Preguntas (el profesor las lee en voz alta con la respuesta):}

\begin{itemize}[leftmargin=1.5cm, itemsep=2pt]
\item[\textcolor{verde}{$\bullet$}] ¿\textcolor{verde}{\textbf{Dónde}} vive Javier? --- \textcolor{verde}{\textbf{En Getafe}}.
\item[\textcolor{colornaranja}{$\bullet$}] ¿\textcolor{colornaranja}{\textbf{Cómo}} se llama su hermana? --- \textcolor{colornaranja}{\textbf{Alejandra}}.
\item[\textcolor{rojo}{$\bullet$}] ¿\textcolor{rojo}{\textbf{Quién}} es Manolo? --- \textcolor{rojo}{\textbf{El abuelo de Javier}}.
\end{itemize}

\vspace{0.3cm}

%% --- TEXTO 2: Lucía ---
\subsection*{\textcolor{naranja}{Texto 2 --- Lucía} {\small (visible en pantalla)}}

{\small
``Lucía Alonso es de un pueblo de Cantabria, España. Tiene \textcolor{morado}{\textbf{tres}} hermanos. Ana, la hermana mayor, estudia \textcolor{azul}{\textbf{ingeniería}}; Carlos, que tiene 16 años, estudia en el mismo instituto que Lucía; y Pablo, que es el hermano pequeño, tiene 9 años. El padre de Lucía es \textcolor{azul}{\textbf{agricultor}} y también tiene vacas. Su madre trabaja \textcolor{verde}{\textbf{en el hotel del pueblo}}. Lucía estudia el primer curso de Educación Secundaria Obligatoria en un pueblo cerca de su casa y también juega fútbol y estudia piano.''}

\vspace{0.2cm}

\textit{Preguntas (el profesor las lee en voz alta con la respuesta):}

\begin{itemize}[leftmargin=1.5cm, itemsep=2pt]
\item[\textcolor{morado}{$\bullet$}] ¿\textcolor{morado}{\textbf{Cuántos}} hermanos tiene Lucía? --- \textcolor{morado}{\textbf{Tres}}.
\item[\textcolor{azul}{$\bullet$}] ¿\textcolor{azul}{\textbf{Qué}} estudia Ana? --- \textcolor{azul}{\textbf{Ingeniería}}.
\item[\textcolor{azul}{$\bullet$}] ¿\textcolor{azul}{\textbf{Qué}} profesión tiene el padre? --- \textcolor{azul}{\textbf{Agricultor}}.
\item[\textcolor{verde}{$\bullet$}] ¿\textcolor{verde}{\textbf{Dónde}} trabaja su madre? --- \textcolor{verde}{\textbf{En el hotel del pueblo}}.
\end{itemize}

\vspace{0.3cm}

\begin{tcolorbox}[colback=gris, colframe=grismedio, boxrule=0.5pt, arc=2pt]
\textit{Instrucción para el profesor:} Lee cada pregunta-respuesta señalando el color y el icono. \textbf{No explica nada.} Exposición pura.
\end{tcolorbox}

\newpage

%% ============================================================
%% DIAPOSITIVA 2 — ¿CÓMO PREGUNTAMOS POR ESTA INFORMACIÓN?
%% ============================================================

\diapheader{2}{¿CÓMO PREGUNTAMOS POR ESTA INFORMACIÓN?}

\vspace{0.2cm}

\begin{marsbox}[Fase MARS: A (Awareness --- noticing guiado + procesamiento receptivo)]
\textit{Técnicas: segmentación visual con colores (Romo Simón, 2014), emparejamiento complemento$\to$interrogativo (VanPatten, 1996), pares mínimos por contraste (Conti, 2023). El estudiante ya ha percibido la correspondencia interrogativo$\leftrightarrow$respuesta (diap.~1). Ahora el profesor hace explícita esa correspondencia moviendo complementos coloreados hacia el interrogativo correcto. El concepto (lugar, cantidad, persona, etc.) aparece DESPUÉS como etiqueta. La parte 2 retira el color para verificar adquisición. La parte 3 presenta pares de contraste.}
\end{marsbox}

\vspace{0.3cm}

%% --- PARTE 1 ---
\subsection*{\textcolor{naranja}{Parte 1 --- El profesor mueve los complementos al interrogativo}}

{\small Textos breves con complementos marcados en color. En barra lateral, los interrogativos sin color: \textit{dónde / de dónde --- cómo --- qué --- quién / quiénes --- cuántos/cuántas --- cuál --- cuándo}. El profesor mueve cada complemento al interrogativo correcto. Al conectar, se forma la pregunta y aparece debajo el concepto.}

\vspace{0.2cm}

%% Javier
\noindent\textbf{Javier} tiene \textcolor{morado}{\textbf{12 años}}, estudia en \textcolor{verde}{\textbf{el instituto de su barrio}}, vive con \textcolor{rojo}{\textbf{su abuelo Manolo}} y come \textcolor{verde}{\textbf{en el comedor del instituto}} de lunes a viernes.

\vspace{0.1cm}

\begin{itemize}[leftmargin=1cm, itemsep=1pt]
\item \textcolor{morado}{\textbf{12 años}} $\to$ mueve a \textbf{cuántos} $\to$ \textit{¿Cuántos años tiene?} $\to$ \textsc{cantidad}
\item \textcolor{verde}{\textbf{el instituto de su barrio}} $\to$ mueve a \textbf{dónde} $\to$ \textit{¿Dónde estudia?} $\to$ \textsc{lugar}
\item \textcolor{rojo}{\textbf{su abuelo Manolo}} $\to$ mueve a \textbf{quién} $\to$ \textit{¿Quién vive con Javier?} $\to$ \textsc{persona}
\item \textcolor{verde}{\textbf{en el comedor del instituto}} $\to$ mueve a \textbf{dónde} $\to$ \textit{¿Dónde come?} $\to$ \textsc{lugar}
\end{itemize}

\vspace{0.15cm}

%% Lucía
\noindent\textbf{Lucía} es de \textcolor{verde}{\textbf{un pueblo de Cantabria}}, tiene \textcolor{morado}{\textbf{tres hermanos}}, estudia \textcolor{azul}{\textbf{piano}} y su madre trabaja en \textcolor{verde}{\textbf{el hotel del pueblo}}.

\vspace{0.1cm}

\begin{itemize}[leftmargin=1cm, itemsep=1pt]
\item \textcolor{verde}{\textbf{un pueblo de Cantabria}} $\to$ mueve a \textbf{de dónde} $\to$ \textit{¿De dónde es?} $\to$ \textsc{origen}
\item \textcolor{morado}{\textbf{tres hermanos}} $\to$ mueve a \textbf{cuántos} $\to$ \textit{¿Cuántos hermanos tiene?} $\to$ \textsc{cantidad}
\item \textcolor{azul}{\textbf{piano}} $\to$ mueve a \textbf{qué} $\to$ \textit{¿Qué estudia?} $\to$ \textsc{cosa/actividad}
\item \textcolor{verde}{\textbf{el hotel del pueblo}} $\to$ mueve a \textbf{dónde} $\to$ \textit{¿Dónde trabaja su madre?} $\to$ \textsc{lugar}
\end{itemize}

\vspace{0.15cm}

%% Carlos
\noindent\textbf{Carlos} tiene \textcolor{morado}{\textbf{16 años}}, estudia \textcolor{verde}{\textbf{en el mismo instituto que Lucía}} y su hermana mayor se llama \textcolor{colornaranja}{\textbf{Ana}}.

\vspace{0.1cm}

\begin{itemize}[leftmargin=1cm, itemsep=1pt]
\item \textcolor{morado}{\textbf{16 años}} $\to$ mueve a \textbf{cuántos} $\to$ \textit{¿Cuántos años tiene Carlos?} $\to$ \textsc{cantidad}
\item \textcolor{verde}{\textbf{en el mismo instituto que Lucía}} $\to$ mueve a \textbf{dónde} $\to$ \textit{¿Dónde estudia?} $\to$ \textsc{lugar}
\item \textcolor{colornaranja}{\textbf{Ana}} $\to$ mueve a \textbf{cómo} $\to$ \textit{¿Cómo se llama su hermana mayor?} $\to$ \textsc{nombre}
\end{itemize}

\vspace{0.4cm}

%% --- PARTE 2 ---
\subsection*{\textcolor{naranja}{Parte 2 --- El estudiante prueba sin color}}

{\small Dos frases nuevas sin color ni icono. El estudiante dice a mano alzada qué interrogativo corresponde a cada complemento. Si acierta, el interrogativo toma el color correspondiente.}

\vspace{0.2cm}

\noindent\textbf{Catalina} es enfermera, trabaja en el hospital, tiene dos hijos y su mejor amiga se llama Rosa.

\begin{itemize}[leftmargin=1cm, itemsep=1pt]
\item enfermera $\leftarrow$ ¿\rule{1.5cm}{0.4pt}? $\to$ (\textbf{qué})
\item el hospital $\leftarrow$ ¿\rule{1.5cm}{0.4pt}? $\to$ (\textbf{dónde})
\item dos hijos $\leftarrow$ ¿\rule{1.5cm}{0.4pt}? $\to$ (\textbf{cuántos})
\item Rosa $\leftarrow$ ¿\rule{1.5cm}{0.4pt}? $\to$ (\textbf{cómo})
\end{itemize}

\vspace{0.15cm}

\noindent\textbf{Pablo} es el hermano pequeño de Lucía, tiene 9 años y su mejor amigo se llama David.

\begin{itemize}[leftmargin=1cm, itemsep=1pt]
\item 9 años $\leftarrow$ ¿\rule{1.5cm}{0.4pt}? $\to$ (\textbf{cuántos})
\item David $\leftarrow$ ¿\rule{1.5cm}{0.4pt}? $\to$ (\textbf{cómo})
\end{itemize}

\vspace{0.4cm}

%% --- PARTE 3 ---
\subsection*{\textcolor{naranja}{Parte 3 --- Pares: misma persona, diferente pregunta}}

{\small Tres pares lado a lado. En cada par, la misma persona pero dos preguntas distintas. El estudiante nota que el interrogativo cambia según la información que se busca.}

\vspace{0.2cm}

\begin{center}
\renewcommand{\arraystretch}{1.6}
\begin{tabular}{|p{6.5cm}|p{6.5cm}|}
\hline
\rowcolor{naranja!10}
\multicolumn{2}{|c|}{\textbf{Par 1 --- Javier}} \\
\hline
¿\textbf{Dónde} vive Javier? --- En Getafe. & ¿\textbf{Quién} vive con Javier? --- Su abuelo Manolo. \\
\hline
\rowcolor{naranja!10}
\multicolumn{2}{|c|}{\textbf{Par 2 --- Lucía}} \\
\hline
¿\textbf{Qué} estudia Lucía? --- Piano. & ¿\textbf{Cuántos} hermanos tiene Lucía? --- Tres. \\
\hline
\rowcolor{naranja!10}
\multicolumn{2}{|c|}{\textbf{Par 3 --- Catalina}} \\
\hline
¿\textbf{Dónde} trabaja Catalina? --- En el hospital. & ¿\textbf{Qué} es Catalina? --- Enfermera. \\
\hline
\end{tabular}
\end{center}

\newpage

%% ============================================================
%% DIAPOSITIVA 3 — ¿CUÁNTOS O CUÁNTAS? ¿QUIÉN O QUIÉNES?
%% ============================================================

\diapheader{3}{¿CUÁNTOS O CUÁNTAS? ¿QUIÉN O QUIÉNES?}

\vspace{0.2cm}

\begin{marsbox}[Fase MARS: A (Awareness --- noticing guiado de concordancia + selección receptiva)]
\textit{Técnicas: input contextualizado con imágenes (Efecto de modalidad, Mayer 2009), gramática visual con colores de concordancia (Romo Simón, 2014), noticing guiado con preguntas cerradas (Conti, 2023), selección receptiva a mano alzada (VanPatten, 1996). En las diapositivas 1--2 todos los ejemplos usaron <<cuántos>> (masculino plural) y <<quién>> (singular). Ahora el estudiante descubre que estas formas CAMBIAN: cuánto concuerda en género y número, y quién tiene plural (quiénes). Se usa la familia de David (p.34) como contexto conocido. Colores: \textcolor{azulmasc}{azul} = masculino, \textcolor{rosafem}{rosa} = femenino.}
\end{marsbox}

\vspace{0.3cm}

%% --- PARTE 1 ---
\subsection*{\textcolor{naranja}{Parte 1 --- Input con imágenes y colores de concordancia}}

{\small El profesor lee, el estudiante observa y escucha. Aparecen 6 viñetas con una imagen de referencia (fotos del árbol genealógico de David, p.34). El interrogativo y la palabra que lo determina están en el mismo color: \textcolor{azulmasc}{azul} = masculino, \textcolor{rosafem}{rosa} = femenino.}

\vspace{0.2cm}

\renewcommand{\arraystretch}{1.6}
\begin{tabular}{@{}c p{13cm}@{}}
\textbf{1.} & \textit{(imagen: Nacho y María)} --- ¿\textcolor{azulmasc}{\textbf{Cuántos}} \textcolor{azulmasc}{\textbf{hermanos}} tiene David? --- Dos: Nacho y María. \\[4pt]
\textbf{2.} & \textit{(imagen: Pilar y Carmen)} --- ¿\textcolor{rosafem}{\textbf{Cuántas}} \textcolor{rosafem}{\textbf{tías}} tiene David? --- Dos: Pilar y Carmen. \\[4pt]
\textbf{3.} & \textit{(imagen: Álvaro)} --- ¿\textcolor{azulmasc}{\textbf{Cuántos}} \textcolor{azulmasc}{\textbf{primos}} tiene David? --- Uno: Álvaro. \\[4pt]
\textbf{4.} & \textit{(imagen: Paloma)} --- ¿\textcolor{rosafem}{\textbf{Cuántas}} \textcolor{rosafem}{\textbf{primas}} tiene David? --- Una: Paloma. \\[4pt]
\textbf{5.} & \textit{(imagen: Alicia sola)} --- ¿\textbf{Quién} es la madre de David? --- \textbf{Alicia}. \\[4pt]
\textbf{6.} & \textit{(imagen: Carlos y Juana juntos)} --- ¿\textbf{Quiénes} son los abuelos de David? --- \textbf{Carlos y Juana}. \\
\end{tabular}

\vspace{0.15cm}

{\small\textit{El profesor lee dos veces. La segunda vez, más despacio, señala los colores con el dedo.}}

\vspace{0.4cm}

%% --- PARTE 2 ---
\subsection*{\textcolor{naranja}{Parte 2 --- Noticing guiado}}

{\small Las 6 viñetas siguen visibles. El profesor hace preguntas cerradas:}

\vspace{0.2cm}

\begin{tcolorbox}[colback=naranja!5, colframe=naranja!60, boxrule=0.5pt, arc=2pt, left=6pt, right=6pt]
\begin{itemize}[leftmargin=0.5cm, itemsep=3pt]
\item \textit{``Mirad las preguntas 1 y 2. ¿A qué palabra acompaña <<cuántos>>?''} $\to$ hermanos. \textit{``¿Y <<cuántas>>?''} $\to$ tías.
\item \textit{``¿Qué diferencia hay entre <<hermanos>> y <<tías>>?''} $\to$ El estudiante nota: masculino / femenino.
\item \textit{``Ahora mirad las preguntas 3 y 4. ¿A qué palabra acompaña <<cuántos>>?''} $\to$ primos. \textit{``¿Y <<cuántas>>?''} $\to$ primas.
\item \textit{``¿Cuántos cambia igual que el nombre que acompaña?''} $\to$ Sí: masculino con masculino, femenino con femenino.
\end{itemize}
\end{tcolorbox}

\vspace{0.2cm}

\begin{center}
\fbox{\parbox{12cm}{\centering\large
\textcolor{azulmasc}{\textbf{cuántos}} + palabra masculina \hspace{1cm} \textcolor{rosafem}{\textbf{cuántas}} + palabra femenina
}}
\end{center}

\vspace{0.2cm}

\begin{tcolorbox}[colback=naranja!5, colframe=naranja!60, boxrule=0.5pt, arc=2pt, left=6pt, right=6pt]
\textit{``Ahora mirad las preguntas 5 y 6. ¿Cuándo uso <<quién>>?''} $\to$ una persona. \textit{``¿Y <<quiénes>>?''} $\to$ más de una persona.
\end{tcolorbox}

\vspace{0.2cm}

\begin{center}
\fbox{\parbox{10cm}{\centering\large
\textbf{quién} $\to$ 1 persona \hspace{2cm} \textbf{quiénes} $\to$ 2 o más personas
}}
\end{center}

\vspace{0.4cm}

%% --- PARTE 3 ---
\subsection*{\textcolor{naranja}{Parte 3 --- Selección receptiva a mano alzada}}

{\small Aparecen 6 frases nuevas con dos opciones entre paréntesis. Los estudiantes eligen a mano alzada la forma correcta. El profesor confirma y la opción correcta toma el color correspondiente.}

\vspace{0.2cm}

\renewcommand{\arraystretch}{1.5}
\begin{tabular}{@{}r p{10cm} l@{}}
\textbf{1.} & ¿(Cuántos / \textcolor{rosafem}{\textbf{Cuántas}}) hermana\textbf{s} tiene Lucía? & \textcolor{rosafem}{\small (hermanas = F pl.)} \\
\textbf{2.} & ¿(Cuántas / \textcolor{azulmasc}{\textbf{Cuántos}}) año\textbf{s} tiene Carlos? & \textcolor{azulmasc}{\small (años = M pl.)} \\
\textbf{3.} & ¿(\textbf{Quién} / Quiénes) es el padre de David? & {\small (una persona)} \\
\textbf{4.} & ¿(Cuántos / \textcolor{rosafem}{\textbf{Cuántas}}) naranja\textbf{s} come María? & \textcolor{rosafem}{\small (naranjas = F pl.)} \\
\textbf{5.} & ¿(Quién / \textbf{Quiénes}) son los tíos de David? & {\small (más de una persona)} \\
\textbf{6.} & ¿(\textcolor{azulmasc}{\textbf{Cuántos}} / Cuántas) idioma\textbf{s} hablas? & \textcolor{azulmasc}{\small (idiomas = M pl.)} \\
\end{tabular}

\vspace{0.2cm}

\begin{tcolorbox}[colback=gris, colframe=grismedio, boxrule=0.5pt, arc=2pt]
\textit{Si el estudiante falla, el profesor dice: ``Mirad la palabra que acompaña. ¿Es masculina o femenina? ¿Una persona o más?'' y otro estudiante intenta.}
\end{tcolorbox}

\newpage

%% ============================================================
%% DIAPOSITIVA 4 — ¡GRITAD LA PREGUNTA! + ¿CON TILDE O SIN TILDE?
%% ============================================================

\diapheader{4}{¡GRITAD LA PREGUNTA! + ¿CON TILDE O SIN TILDE?}

\vspace{0.2cm}

\begin{marsbox}[Fase MARS: S (Structured Production) + cierre con reflexión sobre la tilde]
\textit{Técnicas: producción oral coral con ritmo rápido (Conti, 2023), contraste cómico con/sin tilde para anclar la regla ortográfica. Tras 40+ exposiciones receptivas y de noticing (diaps.~1--3), el estudiante está preparado para producir. Producción oral rápida en grupo-clase (baja ansiedad). Al final, el contraste con/sin tilde genera un anclaje memorable por el efecto cómico: frases absurdas sin tilde que muestran que el acento no es decorativo sino que cambia el significado.}
\end{marsbox}

\vspace{0.4cm}

%% --- PARTE 1: ¡GRITAD LA PREGUNTA! ---
\subsection*{\textcolor{naranja}{Parte 1 --- ¡GRITAD LA PREGUNTA!}}

{\small Aparecen respuestas sueltas de personajes conocidos, una a una, con ritmo rápido. Los estudiantes gritan la pregunta completa que produce esa respuesta. El profesor valida con $\checkmark$ o $\times$. No hay apoyo de color ni iconos --- producción libre.}

\vspace{0.3cm}

\renewcommand{\arraystretch}{1.6}
\begin{center}
\begin{tabular}{|r|>{\bfseries}p{5cm}|p{7cm}|c|}
\hline
\rowcolor{naranja!15}
\textbf{N.} & \textbf{Respuesta} & \textbf{Pregunta esperada} & \\
\hline
1 & En Getafe. & \textit{¿Dónde vive Javier?} & $\checkmark$ \\
\hline
2 & Tres hermanos. & \textit{¿Cuántos hermanos tiene Lucía?} & $\checkmark$ \\
\hline
3 & Alejandra. & \textit{¿Cómo se llama la hermana de Javier?} & $\checkmark$ \\
\hline
4 & Ingeniería. & \textit{¿Qué estudia Ana?} & $\checkmark$ \\
\hline
5 & Dos tías. & \textit{¿Cuántas tías tiene David?} & $\checkmark$ \\
\hline
6 & Carlos y Juana. & \textit{¿Quiénes son los abuelos de David?} & $\checkmark$ \\
\hline
7 & De un pueblo de Cantabria. & \textit{¿De dónde es Lucía?} & $\checkmark$ \\
\hline
8 & Alicia. & \textit{¿Quién es la madre de David?} & $\checkmark$ \\
\hline
9 & Agricultor. & \textit{¿Qué profesión tiene el padre de Lucía?} & $\checkmark$ \\
\hline
10 & Una prima. & \textit{¿Cuántas primas tiene David?} & $\checkmark$ \\
\hline
\end{tabular}
\end{center}

\vspace{0.2cm}

\begin{tcolorbox}[colback=gris, colframe=grismedio, boxrule=0.5pt, arc=2pt]
\textit{Ritmo rápido: máximo 30--40 segundos para las 10 respuestas. Si el grupo responde bien, el profesor dice: ``¡Muy bien! Ya sabéis hacer preguntas. Pero falta algo importante\ldots''}
\end{tcolorbox}

\vspace{0.5cm}

%% --- PARTE 2: ¿CON TILDE O SIN TILDE? ---
\subsection*{\textcolor{naranja}{Parte 2 --- ¿CON TILDE O SIN TILDE?}}

{\small Aparecen frases en pantalla SIN tilde. El profesor pregunta: \textit{``¿Esto es una pregunta?''} Los estudiantes responden. Luego aparece la versión CON tilde y se ve el contraste. Las frases sin tilde dicen algo diferente o absurdo.}

\vspace{0.3cm}

\renewcommand{\arraystretch}{1.2}

%% Par 1
\begin{tcolorbox}[colback=white, colframe=naranja, boxrule=0.8pt, arc=3pt, title={\textbf{Par 1}}, fonttitle=\bfseries\small]
\begin{tabular}{@{}p{7cm} p{7cm}@{}}
\textbf{Sin tilde:} & \textbf{Con tilde:} \\[4pt]
{\large ``Como como un cerdo''} & {\large ``¿\textbf{Cómo} como un cerdo?''} \\[4pt]
\textit{$\to$ No es pregunta. Dice que come igual que un cerdo.} & \textit{$\to$ Pregunta de qué manera se come un cerdo.} \\
\end{tabular}
\end{tcolorbox}

\vspace{0.2cm}

%% Par 2
\begin{tcolorbox}[colback=white, colframe=naranja, boxrule=0.8pt, arc=3pt, title={\textbf{Par 2}}, fonttitle=\bfseries\small]
\begin{tabular}{@{}p{7cm} p{7cm}@{}}
\textbf{Sin tilde:} & \textbf{Con tilde:} \\[4pt]
{\large ``Donde vives hace frío''} & {\large ``¿\textbf{Dónde} vives? Hace frío''} \\[4pt]
\textit{$\to$ No es pregunta. Afirmación: el lugar donde vives hace frío.} & \textit{$\to$ Pregunta dónde vives. Diferente.} \\
\end{tabular}
\end{tcolorbox}

\vspace{0.2cm}

%% Par 3
\begin{tcolorbox}[colback=white, colframe=naranja, boxrule=0.8pt, arc=3pt, title={\textbf{Par 3}}, fonttitle=\bfseries\small]
\begin{tabular}{@{}p{7cm} p{7cm}@{}}
\textbf{Sin tilde:} & \textbf{Con tilde:} \\[4pt]
{\large ``Cuando como no hablo''} & {\large ``¿\textbf{Cuándo} como? No hablo''} \\[4pt]
\textit{$\to$ No es pregunta. Mientras come, no habla.} & \textit{$\to$ Pregunta cuándo come. Y además no habla. Absurdo.} \\
\end{tabular}
\end{tcolorbox}

\vspace{0.2cm}

%% Par 4
\begin{tcolorbox}[colback=white, colframe=naranja, boxrule=0.8pt, arc=3pt, title={\textbf{Par 4}}, fonttitle=\bfseries\small]
\begin{tabular}{@{}p{7cm} p{7cm}@{}}
\textbf{Sin tilde:} & \textbf{Con tilde:} \\[4pt]
{\large ``Cuantos años tienes son importantes''} & {\large ``¿\textbf{Cuántos} años tienes?''} \\[4pt]
\textit{$\to$ No es pregunta. Todos tus años son importantes.} & \textit{$\to$ Pregunta normal. Sin tilde no funciona como pregunta.} \\
\end{tabular}
\end{tcolorbox}

\vspace{0.3cm}

\begin{tcolorbox}[colback=naranja!8, colframe=naranja, boxrule=0.8pt, arc=3pt]
\textit{El profesor cierra: ``¿Veis? La tilde transforma la palabra. Sin tilde no hay pregunta. Todas las palabras que habéis usado hoy llevan tilde siempre que son preguntas.''}

\vspace{0.2cm}

$\to$ \textit{``Abrid el libro en la página 36. Mirad el cuadro de interrogativos. ¿Coincide con lo que habéis descubierto?''}
\end{tcolorbox}

\end{document}
